% GENERATED FILE. Edit canonical lessons, then run npm run generate:books.
% canonical-source-hash: fnv1a64:4c555ce1eb2dabc3
% canonical-lessons: TE-C298-moddu, TE-C298-vankara, TE-C298-cadunaina, TE-C298-bolu, TE-C298-masaka

\chapter{Blunt, Crooked, Flat, Hollow, Dim}
\label{ch:te-blunt-crooked-flat-hollow-dim}
\hlchaptermodality{298}

\begin{chapteropening}
\textbf{By the end of this chapter:} \emph{I can say blunt, crooked, flat, hollow and dim.}

Complete the last lesson of chapter 298: I can say blunt, crooked, flat, hollow and dim.
\end{chapteropening}

\section[\texorpdfstring{\te{మొద్దు} (moddu)}{మొద్దు (moddu)}]{\te{మొద్దు} (moddu) — blunt; dull-witted}

\label{lesson:TE-C298-moddu}



\begin{quote}
Before the new one: say the Telugu for naughty, then the Telugu for honesty.
\end{quote}

\subsection*{You'll want to know: \te{మొద్దు}}
\textbf{\te{మొద్దు}} — \emph{moddu} — \textquotedblleft{}blunt; dull-witted\textquotedblright{}.

\begin{grammarlens}[title={What you've built}]
One more word for this chapter.
\end{grammarlens}

\subsection*{Your turn}
\begin{itemize}
\raggedright
  \item \emph{Say it:} \emph{moddu}
  \item \emph{Say it:} \emph{moddu}, once more
  \item \emph{Say it:} say \emph{moddu}
  \item \emph{Recall:} say the Telugu for naughty, then the Telugu for honesty, then say \emph{moddu} again
\end{itemize}

\begin{tcolorbox}[breakable,colback=teal!4,colframe=teal!35!black,title={Before you move on}]
What does \te{మొద్దు} mean? (\textquotedblleft{}Blunt; dull-witted\textquotedblright{}.) Say it once more.
\end{tcolorbox}
\section[\texorpdfstring{\te{వంకర} (vaṅkara)}{వంకర (vaṅkara)}]{\te{వంకర} (vaṅkara) — crooked, bent}

\label{lesson:TE-C298-vankara}



\begin{quote}
Before the new one: say the Telugu for honesty, then the Telugu for blunt.
\end{quote}

\subsection*{You'll want to know: \te{వంకర}}
\textbf{\te{వంకర}} — \emph{vaṅkara} — \textquotedblleft{}crooked, bent\textquotedblright{}.

\begin{grammarlens}[title={What you've built}]
One more word for this chapter.
\end{grammarlens}

\subsection*{Your turn}
\begin{itemize}
\raggedright
  \item \emph{Say it:} \emph{vaṅkara}
  \item \emph{Say it:} \emph{vaṅkara}, once more
  \item \emph{Say it:} say \emph{vaṅkara}
  \item \emph{Recall:} say the Telugu for honesty, then the Telugu for blunt, then say \emph{vaṅkara} again
\end{itemize}

\begin{tcolorbox}[breakable,colback=teal!4,colframe=teal!35!black,title={Before you move on}]
What does \te{వంకర} mean? (\textquotedblleft{}Crooked, bent\textquotedblright{}.) Say it once more.
\end{tcolorbox}
\section[\texorpdfstring{\te{చదునైన} (cadunaina)}{చదునైన (cadunaina)}]{\te{చదునైన} (cadunaina) — flat, level}

\label{lesson:TE-C298-cadunaina}



\begin{quote}
Before the new one: say the Telugu for blunt, then the Telugu for crooked.
\end{quote}

\subsection*{You'll want to know: \te{చదునైన}}
\textbf{\te{చదునైన}} — \emph{cadunaina} — \textquotedblleft{}flat, level\textquotedblright{}.

\begin{grammarlens}[title={What you've built}]
One more word for this chapter.
\end{grammarlens}

\subsection*{Your turn}
\begin{itemize}
\raggedright
  \item \emph{Say it:} \emph{cadunaina}
  \item \emph{Say it:} \emph{cadunaina}, once more
  \item \emph{Say it:} say \emph{cadunaina}
  \item \emph{Recall:} say the Telugu for blunt, then the Telugu for crooked, then say \emph{cadunaina} again
\end{itemize}

\begin{tcolorbox}[breakable,colback=teal!4,colframe=teal!35!black,title={Before you move on}]
What does \te{చదునైన} mean? (\textquotedblleft{}Flat, level\textquotedblright{}.) Say it once more.
\end{tcolorbox}
\section[\texorpdfstring{\te{బోలు} (bōlu)}{బోలు (bōlu)}]{\te{బోలు} (bōlu) — hollow}

\label{lesson:TE-C298-bolu}



\begin{quote}
Before the new one: say the Telugu for crooked, then the Telugu for flat.
\end{quote}

\subsection*{You'll want to know: \te{బోలు}}
\textbf{\te{బోలు}} — \emph{bōlu} — \textquotedblleft{}hollow\textquotedblright{}.

\begin{grammarlens}[title={What you've built}]
One more word for this chapter.
\end{grammarlens}

\subsection*{Your turn}
\begin{itemize}
\raggedright
  \item \emph{Say it:} \emph{bōlu}
  \item \emph{Say it:} \emph{bōlu}, once more
  \item \emph{Say it:} say \emph{bōlu}
  \item \emph{Recall:} say the Telugu for crooked, then the Telugu for flat, then say \emph{bōlu} again
\end{itemize}

\begin{tcolorbox}[breakable,colback=teal!4,colframe=teal!35!black,title={Before you move on}]
What does \te{బోలు} mean? (\textquotedblleft{}Hollow\textquotedblright{}.) Say it once more.
\end{tcolorbox}
\section[\texorpdfstring{\te{మసక} (masaka)}{మసక (masaka)}]{\te{మసక} (masaka) — dim, blurry}

\label{lesson:TE-C298-masaka}



\begin{quote}
Before the new one: say the Telugu for flat, then the Telugu for hollow.
\end{quote}

\subsection*{You'll want to know: \te{మసక}}
\textbf{\te{మసక}} — \emph{masaka} — \textquotedblleft{}dim, blurry\textquotedblright{}.

\begin{grammarlens}[title={What you've built}]
One more word for this chapter.
\end{grammarlens}

\subsection*{Your turn}
\begin{itemize}
\raggedright
  \item \emph{Say it:} \emph{masaka}
  \item \emph{Say it:} \emph{masaka}, once more
  \item \emph{Say it:} say \emph{masaka}
  \item \emph{Recall:} say the Telugu for flat, then the Telugu for hollow, then say \emph{masaka} again
\end{itemize}

\begin{tcolorbox}[breakable,colback=teal!4,colframe=teal!35!black,title={Before you move on}]
What does \te{మసక} mean? (\textquotedblleft{}Dim, blurry\textquotedblright{}.) Say it once more.
\end{tcolorbox}
